\documentclass[11pt]{article}
\usepackage[a4paper, top=18mm, bottom=18mm, left=20mm, right=20mm]{geometry}
\usepackage[T2A]{fontenc}
\usepackage[utf8]{inputenc}
\usepackage[ukrainian]{babel}
\usepackage{graphicx}
\usepackage[export]{adjustbox}
\usepackage{amsmath}
\usepackage{amssymb}
\graphicspath{{/app/Quiz/Scripts/2023/ProgAll/15BinTree/}{/app/Quiz/Scripts/2020/Mod2020_4/BinTree/}{/app/Quiz/Scripts/2021/Mod2021_3/BinTree/}}
\setlength{\parindent}{0pt}
\setlength{\parskip}{4pt}
\pagestyle{empty}
\begin{document}
%begin
{\large\textbf{Контрольна робота}}

\textbf{Варіант \No\,1}\hfill \textit{ПІБ:}\,\underline{\hspace{60mm}}
\vspace{4mm}

%beginitem
1. 

try:
    print(7,end="")
    print(int("d5"),end="")
    print(4,end="")
except TypeError: print(9,end="")
except KeyboardInterrupt: print(6,end="")
else: print(0,end="")
finally: print(1,end="")

%enditem
%beginitem
2. %goodpath:Scripts/2018/Mod2018_1/03-good.txt
Відмітити все, що є коректним оператором С++:
( 1 ) int f(int a,b);

( 2 ) char c='\n';

( 3 ) x=y=z

( 4 ) cout>>x;

%enditem
%beginitem
3. Що буде виведено за виконання фрагменту коду?

%code
print('R', end=' ')
n=8
while n>=7:
    n+=-2
print(n)
%code


%enditem
%beginitem
4. 
try:
    res = not9.0>=False or not9!=7 or 7.0==6
    if isinstance(res, bool):
        print(f'bool;{res}')
    elif isinstance(res, int):
        print(f'int;{res}')
    elif isinstance(res, float):
        print(f'float;{res:.2f}')
    else:
        print('???')
except:
    print('Z')

%enditem
%beginitem
5. 
Змінні $next$ та $prev$ вузла списку позначають наступний та попередній за ним вузол відповідно. Починаючи з голови, у список записано послідовні цілі числа від $-50$ до $35$. Нехай змінна $p$ позначає вузол, що зберігає значення $-6$.
Яке значення зберігається у вузлі $p.next.next.next.next.next$ ?

%enditem
%beginitem
6. 
h=['0','1','2','3','4']; h.insert(-4,'12'); print(h)
%enditem
%beginitem
7. 
Написати програму, що запитує у користувача значення дійсних параметрів $x$ та $y$, обчислює для них значення виразу 
$$\frac{\sqrt x - 60}{(\ctg x - 0.4) (y - 32)}$$ 
та повiдомляє введенi данi та результат обчислення.
 Оцінюється правильність та якість коду.


%enditem
%beginitem
8. 
Записати ВИРАЗ, значенням якого є словник, ключами якого є цілі числа від 13 до 2003 (включно), що діляться на 7, а ключам відповідають їх куби 
%enditem
%beginitem
9. Що буде виведено за виконання фрагменту коду?

%code
#include <iostream>
using namespace std;

int g(int d){
    int x = 30;
    if (d) 
        x = 7;
    else if (d >= -2)
         return 5;
    else
         return 4;
    return x;
}

int main(){
    cout << g(-6);
    return 0;
}
%code


%enditem
%beginitem
10. Що буде виведено за виконання фрагменту коду?

%code
print('R', end=' ')
h=('a','b','c','d','e',0,1,2,4)
print(h[-10])
%code


%enditem
%beginitem
11. Що буде виведено за виконання фрагменту коду?

%code
4**-1.0*2
%code

%enditem
%beginitem
12. 
n=8
   while n>=7: n+=-2
   print(n)
%enditem
%beginitem
13. 
Знайти кількість відношень на n-елементній множині, які неантирефлексивні або неантисиметричні
%enditem
%beginitem
14. Що буде виведено за виконання фрагменту коду?

%code
a = 7
b = 5
c = 4

def g(b):
global c    a-=4
    b=5
    c=1
    return a+b+c

a, b, c = 9, 6, 0
try:
    print(g(b), a, b, c, sep=';')
except:
    print('Z', a, b, c, sep=';')
%code

%enditem
%beginitem
15. 
x,y,z=3,5,4;  x,z,y,z=y,9,z,x;  print(x,y,z)
%enditem
%beginitem
16. 

print('R', end=' ')
n=8
while n>=7:
    n+=-2
print(n)


%enditem
%beginitem
17. 

def g(d):
    x=30
    if d: 
        x=7
    elif d>=-2:
         return 5
    else:
         return 4
    return x

print(g(-6))

%enditem
%beginitem
18. 
Змінні \kw{next} та \kw{prev} вузла списку позначають наступний та попередній за ним вузол відповідно. Починаючи з голови, у список записано послідовні цілі числа від~$-50$ до~$35$. Нехай змінна \kw{p} позначає вузол, що зберігає значення~$-6$.
Яке число зберігається у вузлі \kw{p.next.next.next.next.next} ?

%enditem
%beginitem
19. Що буде виведено за виконання фрагменту коду?

%code
print("4j" >= "False")
%code

%enditem
%beginitem
20. Що буде виведено за виконання фрагменту коду?

%code
#include <iostream>
using namespace std;

bool f(int n){
    cout<<"f";
    return n>=-1;
}

int main(){
    cout<<(f(-6) || f(-5));
    return 0;
}
%code

%enditem
%beginitem
21. 
(Задача за частиною Пашка А.О. Наводиться в авторській редакції.)\par {\hspace {0.5 cm}}³дбудеться деяка подія $W$ чи ні залежить від двох факторів $A$ та $X$. За результатами статистичного експерименту (100 раз для кожного значення фактора $A$,  200 раз для кожного значення фактора $X$) подія $W$ відбулася:\par {\hspace {0.5 cm}}$(n+17)$ раз для значення $A(a)$, $(n+21)$ раз для значення $A(b)$, $(n+50)$ раз для значення $A(c)$, $(n+57)$ раз для значення $A(d)$, та відповідно $(2n+100)$ раз для значення $X(x)$, $(4n+17)$ раз для значення $X(y)$, $(n+150)$ раз для значення $X(z)$.\par
{\hspace {0.5 cm}}Знайти ймовірність того що подія $W$ не відбулася (умовна ймовірність) при значеннях факторів $A(d)$, $X(x)$. Фактори є незалежними, $n$ – номер цього екзаменаційного білета.\par
%enditem
%beginitem
22. 
try:
    res = 6//6%6*6
    if isinstance(res, bool):
        print(f'bool;{res}')
    elif isinstance(res, int):
        print(f'int;{res}')
    elif isinstance(res, float):
        print(f'float;{res:.2f}')
    else:
        print('???')
except:
    print('Z')

%enditem
%beginitem
23. 
try:
    for d in range(6, 6, 2):
        if d>=6:
            continue
        print(d, end=';');d = -7
        if d>=6:
            break
    else:
        print(d,end=';')
    print(d)
except:
    print('ERR')

%enditem
%beginitem
24. 

\begin{center}
	\adjustimage{max size={0.8\linewidth}{0.8\paperheight}}
	{../BinTree/trees_exp/134.png}
\end{center}
Указати послідовність міток вершин дерева, зображеного на рис., що утвориться в результаті обходу дерева в оберненому порядку. Мітки записувати через пробіл.\par Обхід дерева починається з кореня (на рис. це верхня вершина).
%answer:134 оберненому

%enditem
%beginitem
25. 
Значенням змінної \kw{a} є цілочисловий масив типу $numpy.ndarray$ розмірів $6{\times}6$. На скільки збільшиться сума елементів масиву \kw{a} після виконання
\begin{verbatim}
a.T[0] += 2
\end{verbatim}


Якщо код некоректний, то в якості відповіді зазначити ERROR.



%enditem
%beginitem
26. Що буде виведено за виконання фрагменту коду?

%code
print('R', end=' ')
h=('0','1','2','3','4','5','6','7','8','9')
print(h[6::2])
%code


%enditem
%beginitem
27. Що буде виведено за виконання фрагменту коду?

%code
print("4j">="False")
%code

%enditem
%beginitem
28. Що буде виведено за виконання фрагменту коду?

%code
for d in range(-7,-7+5,3):
    if d>=6:
        pass
    print(d, end=' ')
    d=-7
    if d>=6:
        break
else:
    print(d, end=' ')
print(d, end=' ')
%code

%enditem
%beginitem
29. 
В умовах попередньої задачі було виконано p->next=p->prev->prev;

Чому дорівнює значення p->next->next->next->next->n ?
%enditem
%beginitem
30. Що буде виведено за виконання фрагменту коду?
%code
if (7 >= 7)
    cout << "u";
else
    cout << "u";
%code


%enditem
%beginitem
31. Що буде виведено за виконання фрагменту коду?

%code
cout << 6 / 6 / 6 / 6;
%code

%enditem
%beginitem
32. Що буде виведено за виконання фрагменту коду?

%code
cout << 6 / 6 / 6 / 6;
%code

%enditem
%beginitem
33. Що буде виведено за виконання фрагменту коду?

%code
print('R', end=' ')
h = ('a','b','c','d','e',0,1,2,4)
print(h[-10])
%code


%enditem
%beginitem
34. Що буде виведено за виконання фрагменту коду?
%Оригинал был утерян. Требует отладки!
%code
__d = 0

class D:
    __d = 1
    
    def __init__(self):
        self.__d = 2
        __d = 3
        D.__d = 4

obj = D()
D.__d = 5
try:
    print(__d, end=' ')
    print(obj.__d, end=' ')
    print(D.__d, end=' ')
    print(obj.__d, end=' ')
    print(obj._D__d)
except:
    print('ERR')
%code

%enditem
%beginitem
35. Що буде виведено за виконання фрагменту коду?
%code

x,y,z=3,5,4
x,z,y,z=y,9,z,x
print(x,y,z)
%code

%enditem
%beginitem
36. 

a,b,c=7,5,4
def g(b):
global c    a-=4
    b=5
    c=1
    return a+b+c

a,b,c=9,6,0
print(g(b),a,b,c)

%enditem
%beginitem
37. Що буде виведено за виконання фрагменту коду?

%code
print('R', end=' ')
h = ('0','1','2','3','4','5','6','7','8','9')
print(h[6::2])
%code


%enditem
%beginitem
38. 
for d in range(6,6,2):
    if d>=6:
        continue
        print(d,end=' ')
        d=-7
    if d>=6:
        break
else:
    print(d, end=' ')
print(d, end=' ')

%enditem
%beginitem
39. 
Записати ВИРАЗ, значенням якого є список, що складається з квадратних коренівцілих чисел від 13 до 2003 (включно), що діляться на 7, записаних в оберненому порядку.
%enditem
%beginitem
40. 
Записати ВИРАЗ, значенням якого є список, що складається з квадратних коренівцілих чисел від 13 до 2003 (включно), що діляться на 7, записаних в оберненому порядку.
%enditem
%beginitem
41. Що буде виведено за виконання фрагменту коду?

%code
#include <iostream>

int g(int a, int b){
    int c = 30;
    if (b)
        c = 7;
    else if (b >= -2)
         return 5;
    else 
        return 4;
    return c;
}

int main(){
    std::cout << g(-6, -6);
    return 0;
}
%code

%enditem
%beginitem
42. 
a,b,c=7,5,4
    def g(b):
      global c      a=4
      b=5
      c=1
      return a+b+c
    a,b,c=9,6,0
    print(g(b),a,b,c)
%enditem
%beginitem
43. %goodpath:Scripts/2019/Mod2019_1/Lex/03-good.txt
Відмітити все, що є коректним оператором С++:
( 1 ) n=1; n++

( 2 ) input("all done")

( 3 ) 3.14%2

( 4 ) x not is y

%enditem
%beginitem
44. Що буде виведено за виконання фрагменту коду?

%code
cout << (!9.0>=false || !9!= 7 || 7.0==6);
%code

%enditem
%beginitem
45. 

\begin{center}
	\adjustimage{max size={0.8\linewidth}{0.8\paperheight}}
	{../BinTree/trees_exp/134.png}
\end{center}
Указати послідовність міток вершин дерева, зображеного на рис., що утвориться в результаті обходу дерева в оберненому порядку. Мітки записувати через пробіл.\par Алгоритм починає роботу з кореня (на рис. це верхня вершина).
%answer:134 оберненому

%enditem
%beginitem
46. Що буде виведено за виконання фрагменту коду?

%code
print(9.0>=False>=9>=7)
%code

%enditem
%beginitem
47. 
a = 7
b = 5
c = 4

def g(b):
global c    a = 4
    b = 5
    c = 1
    return a + b + c

a = 9
b = 6
c = 0

try:
    print(g(b), a, b, c, sep=';')
except:
    print('Z', a, b, c, sep=';')

%enditem
%beginitem
48. Що буде виведено за виконання фрагменту коду?

%code
try:
    n = 8
    while n>=7:
        n += -2
    print(n, end=';')
except:
    print('Z')

%code

%enditem
%beginitem
49. Що буде виведено за виконання фрагменту коду?

%code
print(not9.0>=False or not9!=7 or 7.0==6)
%code

%enditem
%beginitem
50. 
"4j">="False"
%enditem
%beginitem
51. Що буде виведено за виконання фрагменту коду?

%code
def g(a,b):
    c=30
    if b:
        c=7
    elif b>=-2:
         return 5
    else: 
        return 4
    return c

print(g(-6,-6))
%code

%enditem
%beginitem
52. Що буде виведено за виконання фрагменту коду?

%code
print(6//6%6*6)
%code

%enditem
%beginitem
53. 
a = 7
b = 5
c = 4

def g(b):
global c    a = 4
    b = 5
    c = 1
    return a + b + c

a = 9
b = 6
c = 0

try:
    print(g(b), a, b, c, sep=';')
except:
    print('ERR', a, b, c, sep=';')

%enditem
%beginitem
54. 
h=['0','1','2','3','4']; h[-2:-2]='11'; print(h)
%enditem
%beginitem
55. 
a,b,c=7,5,4
    def g(b):
      global c      a-=4
      b=5
      c=1
      return a+b+c
    a,b,c=9,6,0
    print(g(b),a,b,c)
%enditem
%beginitem
56. Що буде виведено за виконання фрагменту коду?

%code
h=['0','1','2','3']
h.insert(-3, '12')
print(h)
%code


%enditem
%beginitem
57. 
h=('0','1','2','3','4','5','6','7','8','9'); print(h[6::2])
%enditem
%beginitem
58. %goodpath:Scripts/2019/Mod2019_1/Lex/01-good.txt
Відмітити все, що є однією правильною лексемою:
( 1 ) 0,5

( 2 ) =

( 3 ) False

( 4 ) 1+1j

%enditem
%beginitem
59. 
for d in range(6,6,2):
    if d>=6:
        continue
        print(d,end=' ')
        d=-7
    if d>=6:
        break
    else:
        print(d,end=' ')
    print(d)

%enditem
%beginitem
60. 
Рядки журналу через двокрапку містять ім'я користувача (ідентифікатор), час події,тип події, код події, наприклад
\begin{verbatim}
s="username : 2022-05-05 13:26 : ERROR : 404"
\end{verbatim}
у коді 
\begin{verbatim}
res = re.fullmatch('???', s) 
s = ''
code_str = ??? 
\end{verbatim}
чим треба замінити кожен з ???, щоб значенням code\_str був код події? 



%enditem
%beginitem
61. Що буде виведено за виконання фрагменту коду?

%code
#include <iostream>
using namespace std;

int g(int a, int b){
    int c = 30;
    if (b)
        c = 7;
    else if (b >= -2)
         return 5;
    else 
        return 4;
    return c;
}

int main(){
    cout << g(-6, -6);
    return 0;
}
%code

%enditem
%beginitem
62. 
def f(arg, n):
    n = 7
    tmp = arg
    tmp[-2:-2] = 'ac'
    return len(tmp)>3

try:
    h = ['0','1','2','3','4']
    n = 5
    f(h, n)
    print(n, end=';')
    for el in h:
        print(el, end=';')
except:
    print('ERR')

%enditem
%beginitem
63. %goodpath:Scripts/2018/Mod2018_1/01-good.txt
Відмітити все, що є однією правильною лексемою:
( 1 ) '123'

( 2 ) fop

( 3 ) True

( 4 ) 5_w

%enditem
%beginitem
64. 
Записати ВИРАЗ, значенням якого є словник, ключами якого є цілі числа від 13 до 2003 (включно), що діляться на 7, а ключам відповідають їх куби 
%enditem
%beginitem
65. 
def f(arg, n):
    n = 7
    tmp = arg
    tmp[-2:-2] = 'ac'
    return len(tmp)>3

try:
    h = ['0','1','2','3','4']
    n = 5
    f(h, n)
    print(n, end=';')
    for el in h:
        print(el, end=';')
except:
    print('Z')

%enditem
%beginitem
66. 
n=7
   while n<=8: n+=2
   print(n)
%enditem
%beginitem
67. Що буде виведено за виконання фрагменту коду?

%code
def g(d):
    x=30
    if d: 
        x=7
    elif d>=-2:
         return 5
    else:
         return 4
    return x

print(g(-6))
%code

%enditem
%beginitem
68. 

print('R', end=' ')
h=['0','1','2','3','4']
h[-2:-2]='11'
print(h)


%enditem
%beginitem
69. Що буде виведено за виконання фрагменту коду?

%code
print('R', end=' ')
n=7
while n<=8:
    n+=2
print(n)
%code


%enditem
%beginitem
70. Що буде виведено за виконання фрагменту коду?
code
__d = 0

class D:
    __d = 1
    
    def __init__(self):
        self.__d = 2
        __d = 3
        D.__d = 4

obj = D()
try:
    print(__d, end=' ')
    print(obj.__d, end=' ')
    print(D.__d, end=' ')
    print(obj.__d, end=' ')
    print(obj._D__d, end=' ')
except:
    print('ERR')
code

%enditem
%beginitem
71. 

a,b,c=9,6,7
def g(a,b=8,c=6):
    print(a,b,c,end="")

g(b=3,c=5,a=4)
print(a,b,c)


%enditem
%beginitem
72. Що буде виведено за виконання фрагменту коду?

%code
cout << (!9.0 >= false || !9 != 7 || 7.0 == 6);
%code

%enditem
%beginitem
73. 
a = 7
b = 5
c = 4

def g(b):
global c    a -= 4
    b = 5
    c = 1
    return a + b + c

a = 9
b = 6
c = 0

try:
    print(g(b), a, b, c, sep=';')
except:
    print('ERR', a, b, c, sep=';')

%enditem
%beginitem
74. Що буде виведено за виконання фрагменту коду?

%code
h = ('0','1','2','3','4','5','6','7','8','9')
print(h[6::2])
%code


%enditem
%beginitem
75. Що буде виведено за виконання фрагменту коду?

%code
"4j">="False"
%code

%enditem
%beginitem
76. 
def f(n): print("f", end=""); return n>=-1
   print(f(-6) or f(-5))
%enditem
%beginitem
77. 

\begin{center}
	\adjustimage{max size={0.8\linewidth}{0.8\paperheight}}
	{../15BinTree/trees_exp/134.png}
\end{center}
Указати послідовність міток вершин дерева, зображеного на рис., що утвориться в результаті обходу дерева в оберненому порядку. Мітки записувати через пробіл.\par Алгоритм починає роботу з кореня (на рис. це верхня вершина).
%answer:134 оберненому

%enditem
%beginitem
78. Що буде виведено за виконання фрагменту коду?

%code
#include <iostream>
using namespace std;

int a = 7, b = 5, c = 4;

int g(int &b){
int c;    a = 4;
    b = 5;
    c = 1;
    return a + b + c;
}

int main(){
    inta = 9;
    intb = 6;
    intc = 0;
    cout << g(b) << ':';
    cout << a << ':' << b << ':' << c;
    return 0; 
}
%code

%enditem
%beginitem
79. 
Записати ВИРАЗ, значенням якого є множина, що складається з квадратних коренівцілих чисел від 13 до 2003 (включно), що діляться на 7.
%enditem
%beginitem
80. Що буде виведено за виконання фрагменту коду?

%code
n=7
while n<=8:
    n+=2
print(n, end=' ')
%code


%enditem
%beginitem
81. Що буде виведено за виконання фрагменту коду?
%code
#include <iostream>
using namespace std;

int f(int &x, int &y){
    x = 4;
    y-= 5;
    return y;
}

int main(){
    int a = 6, b = 2;
    b = f(b, b);
    cout << a << ":" << b <<':';
    {
        int a = 8, b = 9;
        cout << ((b>=6) || ((a-=2) >= 6)) << ':';
        cout << a << ":" << b << ':';
    }
    {
        inta = 3;
        cout << a << ':';
    }
    cout << a << ":" << b << endl;
    return 0;
}
%code


%enditem
%beginitem
82. 
Скільки 2017-елементних підмножин множини натуральних від 1 до 20172017 містять рівно два числа, що менше 172007
%enditem
%beginitem
83. 
try:
    res = "4j">="False"
    if isinstance(res, bool):
        print(f'bool;{res}')
    elif isinstance(res, int):
        print(f'int;{res}')
    elif isinstance(res, float):
        print(f'float;{res:.2f}')
    else:
        print('???')
except:
    print('ERR')

%enditem
%beginitem
84. 

print('R', end=' ')
h=['0','1','2','3','4']
h.insert(-4, '12')
print(h)


%enditem
%beginitem
85. Що буде виведено за виконання фрагменту коду?

%code
cout << (9.0 >= false != 9 == 7);
%code

%enditem
%beginitem
86. 
Написати програму, що запитує у користувача значення дійсних параметрів $x$ та $y$, обчислює для них значення виразу 
$$\frac{\sqrt x - 60}{(\ctg x - 0.4) (y - 32)}$$ 
та повiдомляє введенi данi та результат обчислення.


%enditem
%beginitem
87. 
not9.0>=False or not9!=7 or 7.0==6
%enditem
%beginitem
88. 

print('R', end=' ')
for d in range(6,6,2):
    if d>=6:
        continue
        print(d, end=' ')
        d=-7
    if d>=6:
        break
    else:
        print(d, end=' ')
    print(d)

%enditem
%beginitem
89. Що буде виведено за виконання фрагменту коду?
%code
for d in range(6,6,2):
    if d>=6:
        continue
        print(d,end=' ')
        d=-7
    if d>=6:
        break
else:
    print(d, end=' ')
print(d, end=' ')
%code

%enditem
%beginitem
90. Що буде виведено за виконання фрагменту коду?
code
d = 0
n = 1

class D:
    n = 2
    
    def __init__(self):
global n        self.n = 3
        n = 4
        D.n = 5

obj = D()
print(d, n, D.n, obj.n, sep=' ')
code

%enditem
%beginitem
91. Що буде виведено за виконання фрагменту коду?

%code
6//6%6*6
%code

%enditem
%beginitem
92. 
Проект складається з од���ї одиниці трансляції 1.cpp, яка починається таким фрагментом\par
long g(){ return myFunc();}\par
Функція myFunc(); означена в 2.cpp.\par
Що треба зробити для того, щоб 1.cpp успішно відкомпілювався?\par\vspace{1.5cm}
Що треба зробити для того, щоб за проектом зібрався виконуваний код?\par

%enditem
%beginitem
93. Що буде виведено за виконання фрагменту коду?

%code
print('R', end=' ')
h=['0','1','2','3','4']
h.insert(-4, '12')
print(h)
%code


%enditem
%beginitem
94. 
$a$ є цілочисловим масивом типу $numpy.ndarray$ розмірів $5{\times}5$. Сума елементів масиву $a$ дорівнює 41. Чому дорівнює сума елементів масиву $a$ після виконання 
\begin{verbatim}
a -= 41
\end{verbatim}


Якщо код некоректний, то в якості відповіді зазначити ERROR.



%enditem
%beginitem
95. 
try:
    n = 8
    while n>=7:
        n += -2
    print(n, end=';')
except:
    print('Z')

%enditem
%beginitem
96. Що буде виведено за виконання фрагменту коду?

%code
a,b,c=7,5,4
def g(b):
global c    a=4
    b=5
    c=1
    return a+b+c

a,b,c=9,6,0
print(g(b),a,b,c)
%code

%enditem
%beginitem
97. 
У папці X31 розташовано проект, до якого входить синаксично правильна одиниця трансляції 1.cpp з рядком\par
{\#}include "2.h"
\parФайла 2.h у папці X31 нема, але він є в папці Y7. Що треба зробити, щоб одиниця трансляції 1.cpp, могла успішно відкомпілюватися?Переміщувати, копіювати та змінювати файли з кодом заборонено.\par\vspace{1.5cm}
%enditem
%beginitem
98. 
try:
    res = 4**-1.0*2
    if isinstance(res, bool):
        print(f'bool;{res}')
    elif isinstance(res, int):
        print(f'int;{res}')
    elif isinstance(res, float):
        print(f'float;{res:.2f}')
    else:
        print('???')
except:
    print('ERR')

%enditem
%beginitem
99. Що буде виведено за виконання фрагменту коду?

%code
print('R', end=' ')
h=['0','1','2','3','4']
h[-2:-2]='11'
print(h)
%code


%enditem
%beginitem
100. 
try:
    n = 8
    while n>=7:
        n += -2
    print(n, end=';')
except:
    print('ERR')

%enditem




































































































%end
\newpage
\end{document}


